\documentclass[11pt,a4paper]{article}
\usepackage[margin=2.1cm]{geometry}
\usepackage[T1]{fontenc}
\usepackage{lmodern}
\usepackage{amsmath,amssymb,booktabs,tabularx}
\usepackage{xcolor}
\usepackage{microtype}
\usepackage{fancyhdr}
\usepackage[colorlinks=true,urlcolor=blue!55!black,linkcolor=blue!55!black,citecolor=blue!55!black]{hyperref}
\definecolor{ink}{HTML}{163B50}
\setlength{\parindent}{0pt}
\setlength{\parskip}{6pt}
\setlength{\headheight}{14pt}
\pagestyle{fancy}
\fancyhf{}
\fancyhead[L]{\small\textcolor{ink}{Reionization: ionization-rate benchmarks}}
\fancyhead[R]{\small $z\simeq5.5$}
\fancyfoot[C]{\small\thepage}
\newcommand{\Gyr}{\mathrm{Gyr}}
\newcommand{\HI}{\mathrm{HI}}
\newcommand{\HII}{\mathrm{HII}}
\begin{document}
\thispagestyle{plain}
{\LARGE\bfseries\textcolor{ink}{Ionization rates near the end\\[3pt] of reionization}}\par
{\large Global budget per Hubble time and photoionization rate at $z=5.5$}

At $z=5.5$, the observational benchmarks considered here are approximately
\textbf{5 ionizing photons per hydrogen nucleus per Hubble time} for the
global photon supply, and $\boldsymbol{\Gamma_{\HI}\simeq
3.4\times10^{-13}\,\mathrm{s}^{-1}}$ per neutral hydrogen atom.
The conventional recombination-maintenance requirement is lower.
These values retain the cited papers' assumptions; they do not select a
reionization history for the research project.

\section*{1. Global budget per hydrogen nucleus per Hubble time}
Define the Hubble time at the epoch of interest and the dimensionless
photon-supply rate as
\begin{equation}
 t_H(z)\equiv\frac{1}{H(z)},\qquad
 \mathcal N_H(z)\equiv
 \frac{\dot n_{\rm ion}(z)}{\bar n_{\rm H}}\,t_H(z).
\end{equation}
Here $\dot n_{\rm ion}$ is the ionizing-photon emissivity escaping into the
intergalactic medium (IGM), and $\bar n_{\rm H}$ counts \emph{all hydrogen
nuclei}; both densities use consistent comoving units.

\textbf{Observationally inferred supply.}
Gaikwad et al.\ \cite{Gaikwad2023}, Table 3, report
\begin{equation}
 \dot n_{\rm ion}(5.5)=
 (0.587^{+0.417}_{-0.149})\times10^{51}
 \,\mathrm{photons\,s^{-1}\,cMpc^{-3}}.
\end{equation}
Using their cosmology ($h=0.678$, $\Omega_m=0.308$,
$\Omega_\Lambda=0.692$, $\Omega_b=0.0482$, and helium mass fraction
$Y=0.24$) gives $t_H(5.5)=1.56\,\Gyr$ and
\begin{equation}
 \boxed{\mathcal N_H(5.5)=5.2^{+3.7}_{-1.3}
 \quad\text{photons per H nucleus per Hubble time}.}
\end{equation}
The $1\sigma$ uncertainties are propagated from the published emissivity.
The inference uses an absorption-limited radiation field, source spectral
index $\alpha_s=2.0\pm0.6$, and absorber-distribution slope
$\beta_{\HI}=1.3\pm0.05$ \cite[\S6.4]{Gaikwad2023}.
Converting the escaping photons into actual hydrogen ionizations requires
accounting for their absorption and any secondary ionizations.

\textbf{Conventional maintenance requirement.}
A commonly used prescription assumes $C_{\HII}=3$,
$T=20{,}000\,\mathrm K$, case-B recombination, and an electron contribution
from singly ionized helium \cite{Robertson2015}. The numerical approximation
\cite[Eq.\ (3)]{Lovell2018}
\begin{equation}
 t_{\rm rec}\simeq0.93\,\Gyr
 \left(\frac{C_{\HII}}{3}\right)^{-1}
 \left(\frac{T}{2\times10^4\,\mathrm K}\right)^{0.7}
 \left(\frac{1+z}{7}\right)^{-3}
\end{equation}
gives $t_{\rm rec}(5.5)\simeq1.16\,\Gyr$. Using Robertson et al.'s
cosmology to evaluate $t_H$ therefore yields
\begin{equation}
 \boxed{\mathcal N_{H,\rm crit}(5.5)=\frac{t_H}{t_{\rm rec}}
 \simeq1.3
 \quad\text{ionizations per H nucleus per Hubble time}.}
\end{equation}
This replaces recombinations in an already ionized IGM. Higher inferred
clumping factors increase the requirement; Davies et al.\ \cite{Davies2024}
infer $C\approx12$, normalized using case-B recombination at $10^4\,\mathrm K$.

\textbf{Meaning of the time unit.} ``Per Hubble time'' is an instantaneous
rate normalization, not a photon count accumulated over a Hubble time of
evolving emissivity. Here $t_H$ means $1/H(z)$, not the present-day Hubble
time or the age of the Universe.

\newpage
\section*{2. Photoionization rate per neutral hydrogen atom}
At the same redshift, Gaikwad et al.\ \cite[Table 3]{Gaikwad2023} report
\begin{equation}
 \boxed{\Gamma_{\HI}(5.5)=
 \left(3.44^{+2.19}_{-1.30}\right)\times10^{-13}
 \,\mathrm{s}^{-1}.}
\end{equation}
The central value corresponds to a photoionization timescale
\begin{equation}
 \Gamma_{\HI}^{-1}\simeq92{,}000\,\mathrm{yr}
\end{equation}
for a neutral atom exposed to that background. The rate is inferred from
Ly$\alpha$-forest observations through the authors' simulation modeling;
the quoted uncertainties are $1\sigma$.

The normalization differs from the global budget because $\Gamma_{\HI}$
acts on \textbf{neutral atoms}. Locally, the photoionization-event rate per
total hydrogen nucleus is $x_{\HI}\Gamma_{\HI}$, where $x_{\HI}$ is the
neutral fraction. A global value requires the appropriate spatial average;
one cannot generally replace this average by the product of separately
averaged quantities.

\section*{What these benchmarks establish}
\renewcommand{\arraystretch}{1.25}
\begin{tabularx}{\textwidth}{@{}X l@{}}
\toprule
\textbf{Quantity at $z=5.5$} & \textbf{Benchmark} \\
\midrule
Inferred escaping-photon supply per H nucleus per $t_H$
& $5.2^{+3.7}_{-1.3}$ \\
Conventional maintenance ionizations per H nucleus per $t_H$
& $\simeq1.3$ \\
Photoionization rate per neutral H atom
& $(3.44^{+2.19}_{-1.30})\times10^{-13}\,\mathrm{s}^{-1}$ \\
\bottomrule
\end{tabularx}

\textbf{Completion by $z\simeq5.5$ still requires an ionization history.}
In the conventional photon-budget model \cite{Robertson2015,Davies2024},
\begin{equation}
 \frac{dQ_{\HII}}{dt}=
 \frac{\dot n_{\rm ion}}{\bar n_{\rm H}}-
 \frac{Q_{\HII}}{t_{\rm rec}}.
\end{equation}
The maintenance benchmark follows by setting $Q_{\HII}=1$ and
$dQ_{\HII}/dt=0$. Reaching that state also depends on the earlier photon
supply, initial ionized fraction, and recombination history. Neither the
maintenance benchmark nor the measured $\Gamma_{\HI}$ alone establishes
that enough earlier ionizations occurred.

\begin{thebibliography}{9}
\small
\setlength{\itemsep}{3pt}
\bibitem{Gaikwad2023}
P. Gaikwad et al.\ (2023).
\emph{Measuring the photo-ionization rate, neutral fraction and mean free
path of HI ionizing photons at $4.9\leq z\leq6.0$ from a large sample of
XShooter and ESI spectra}. Table 3 and \S6.4.
\href{https://arxiv.org/abs/2304.02038}{arXiv:2304.02038}.

\bibitem{Robertson2015}
B. E. Robertson et al.\ (2015).
\emph{Cosmic Reionization and Early Star-Forming Galaxies:
A Joint Analysis of New Constraints from Planck and the Hubble Space Telescope}.
Eq.\ (4) and adopted recombination prescription.
\href{https://arxiv.org/abs/1502.02024}{arXiv:1502.02024}.

\bibitem{Lovell2018}
M. R. Lovell et al.\ (2018).
\emph{ETHOS---an effective theory of structure formation: predictions for
the high-redshift Universe---abundance of galaxies and reionization}.
Eq.\ (3).
\href{https://arxiv.org/abs/1711.10497}{arXiv:1711.10497}.

\bibitem{Davies2024}
F. B. Davies, S. E. I. Bosman, and S. R. Furlanetto (2024).
\emph{The Predicament of Absorption-dominated Reionization II:
Observational Estimate of the Clumping Factor at the End of Reionization}.
\href{https://arxiv.org/abs/2406.18186}{arXiv:2406.18186}.
\end{thebibliography}
\end{document}
